\chapter{Design}
\label{ch:design}

This chapter describes \textit{what} the system does and \textit{why} each
high-level decision was taken.
The companion chapter on implementation (Chapter~\ref{ch:impl}) describes
\textit{how} the design is realised in code: the specific libraries, file
layout, and engineering choices.
The separation follows the principle of presenting the conceptual
architecture first---in the vocabulary of music notation and machine
learning---so that the reasoning behind each choice is legible before the
reader is asked to interpret library names and configuration flags.

% ─────────────────────────────────────────────────────────────────────────────
\section{Problem Decomposition}
\label{sec:design-pipeline}

The end-to-end transcription task is decomposed into five sequential
stages, with a clearly defined contract between each pair.
Figure~\ref{fig:pipeline-design} summarises the data flow.

\begin{enumerate}
    \item \textbf{Preprocessing} takes the input bytes (a PDF page or a
          raster image) and returns a clean, deskewed binary image of
          the page.
          This stage is purely image-processing; it is independent of any
          musical content.
    \item \textbf{Staff detection} takes the binary page and returns a
          list of detected \term{systems}, each consisting of a cropped
          staff image, the strip of pixels above it (where chord symbols
          live), and the bounding box on the page.
          This is the only stage that performs spatial layout analysis;
          everything downstream operates on one staff at a time.
    \item \textbf{Melody recognition} takes a single staff image and
          returns a flat sequence of LMX tokens describing the notes,
          rests, key, time signature, and barlines printed on that
          staff.
          This is the core machine-learning stage.
    \item \textbf{Grammar correction} takes the raw token sequence and
          returns a syntactically valid LMX sequence, repairing
          structural errors (misplaced accidentals, malformed time
          signatures, runaway measures) without inventing new musical
          content.
          This stage is deterministic, not learned.
    \item \textbf{Chord recognition} takes the chord-strip image
          associated with each staff and returns a list of jazz chord
          symbols.
          It is independent of the melody recognition stage and uses a
          separate neural network.
\end{enumerate}

\begin{figure}[htbp]
\centering
\todo{Block diagram: PDF/image → Preprocess → Staff Detect → \{Melody
CRNN, Chord CRNN\} → Grammar Fix + Chord Postprocess → JSON.
Label each arrow with the data type that flows along it.}
\caption{High-level data flow of the five-stage inference pipeline.
Each block has a single responsibility and a typed input/output
contract.}
\label{fig:pipeline-design}
\end{figure}

\textbf{Why a five-stage decomposition.}
Three constraints drove this decomposition.
First, the chord-symbol stream and the melody stream require
fundamentally different recognisers---one is staff-relative pattern
recognition, the other is text-like OCR---so they must be separable.
Second, the grammar of the LMX output is more easily enforced by a
deterministic validator than learned end-to-end from a finite training
set; isolating that logic in a dedicated stage produces more reliable
output without enlarging the model.
Third, both major surveys of the field~\cite{calvoZaragozaUnderstandingOMR2020,shatriOPTICALMUSICRECOGNITION}
explicitly recommend modular decomposition for OMR systems that target
specific notation styles, because it allows each module to be replaced
or retrained independently as the target domain evolves.

% ─────────────────────────────────────────────────────────────────────────────
\section{Data Strategy}
\label{sec:design-data}

The absence of a labelled corpus of jazz lead sheets is the single most
important constraint on this project.
The data strategy is therefore designed around two complementary ideas:
\textit{synthetic supervision} for breadth and \textit{domain
augmentation} for realism.

\subsection{Synthetic supervision from PrIMuS}

PrIMuS~\cite{PrIMuSDataset} provides 87\,678 monophonic music incipits,
each accompanied by ground-truth symbolic annotations (clef, key, time
signature, and the sequence of notes and rests on the staff).
The corpus is, however, distributed with renderings in the Sibelius
engraving style, which does not match the Real Book aesthetic that the
system must face at inference time.
The chosen strategy is to keep the \textit{semantic content} of every
PrIMuS sample but to discard the supplied renderings and re-render each
sample in the LilyJAZZ engraving style.
This yields a corpus that is large enough to train a deep model, free of
manual annotation cost, and visually aligned with the target domain.

Per-sample rendering variations sample the natural visual variability of
real lead sheets without altering the musical content: the time
signature is rendered numerically or as the C / cut-time glyph with
equal probability; cautionary accidentals are added in roughly half of
samples; the global staff size is drawn uniformly from a six-element
range, varying note size and horizontal density;
and unbeamed eighth notes appear in a small fraction of samples to
expose the model to both beamed and unbeamed groupings.

\subsection{Domain augmentation: from clean renders to scanned pages}

Even after the engraving style is corrected, a model trained only on
clean LilyJAZZ renders generalises poorly to real scans: the feature
distribution shifts substantially when crisp black symbols on a white
background are replaced by blurred, uneven, noisy images with variable
contrast and geometric distortion.
The system therefore generates a second, augmented version of every
training image that simulates the distortions characteristic of
real-world scans: page warping and lens distortion, scan misalignment,
focus blur, exposure variation, sensor noise, uneven illumination, ink
bleed, halftone banding, and JPEG compression.
The model is trained on the union of clean and augmented samples, with
the augmented copies sharing the same ground-truth label.
This is the same recipe followed by every recent end-to-end OMR
system~\cite{calvoZaragozaCameraPrimus2018,mayerPracticalEndToEnd2024,riosVilaSMTPlusPlus2024,yangLEGATO2025};
the specific list of distortions used here is calibrated against a
sample of real Real Book scans.

\subsection{Domain restriction: filtering the training set}

PrIMuS contains many samples that lie outside the lead-sheet domain:
multi-staff renders, non-treble clefs, exotic time signatures.
Three deliberate filters are applied at dataset construction time to
restrict training to the in-scope distribution: multi-staff samples are
dropped, samples with non-treble clefs are dropped, and samples with
time signatures outside the eight-element jazz common-time set
($4/4$, $3/4$, $2/4$, $2/2$, $6/8$, $6/4$, $5/4$, $12/8$) are dropped.
This is a deliberate \textit{contract}, not a workaround: disabling the
filters would force the model to spend capacity on visual patterns that
never appear at inference, without improving its performance on the
target distribution.
The same contract is enforced by the grammar fixer at inference time,
which rewrites any predicted clef other than G2 to G2 and rejects
time signatures outside the allowed set.

\subsection{Fine-tuning on real pages}

The synthetic-to-real gap cannot be fully closed by augmentation alone:
the visual distribution of real Real Book scans has detail---specific
ink-pen artefacts, paper grain, page-layout idiosyncrasies---that no
practical augmentation pipeline can reproduce exactly.
The design therefore plans for a second training phase: after the
synthetic model has converged, a small set of manually annotated Real
Book pages is injected into the training split (only) and the model is
fine-tuned on the combined corpus.
The fine-tuning stage is intentionally lightweight; its purpose is to
adjust the feature distribution rather than to teach the model new
musical concepts.

% ─────────────────────────────────────────────────────────────────────────────
\section{Output Representation: LMX}
\label{sec:lmx-format}

The model's output is a flat sequence of LMX tokens.
This section defines the LMX grammar used in this thesis, describes the
vocabulary the model learns, and justifies the design choices made.

\subsection{Grammar}

A valid LMX sequence is a sequence of measures, each introduced by a
\code{measure} token and optionally followed by a clef, key, or time
declaration before its musical content.
Within a measure, the content is a flat sequence of notes and rests; no
nesting or polyphony is permitted.

\begin{quote}\small
\begin{verbatim}
SEQUENCE  := HEADER ELEMENT* (BARLINE ELEMENT*)*
HEADER    := "measure" [CLEF] [KEY] [TIME]
BARLINE   := "measure" [KEY] [TIME]
CLEF      := "clef:G2"
KEY       := "key:fifths:" <fifths>
TIME      := "time" "beats:" <n> "beat-type:" <d>
ELEMENT   := NOTE | REST
NOTE      := "pitch:" <A-G> "octave:" <0-8> DURATION DOT*
             [ACCIDENTAL] [TIE_STOP] [FERMATA]
REST      := "rest" DURATION DOT* [FERMATA]
DURATION  := "whole"|"half"|"quarter"|"eighth"|"16th"|"32nd"|"64th"
ACCIDENTAL:= "flat" | "sharp" | "natural"
TIE       := "tied:start" ... "tied:stop"
\end{verbatim}
\end{quote}

Three encoding choices in this grammar deserve specific justification.

\textbf{Three-token time signatures.}
The time signature is emitted as three separate tokens
(\code{time}, \code{beats:n}, \code{beat-type:d}) rather than a single
joint token \code{time:n/d}.
This factorisation reduces the vocabulary size from forty-plus possible
combinations to roughly a dozen \code{beats:} and \code{beat-type:}
values, and it lets the model share representations between related
time signatures ($4/4$ and $2/4$ share \code{beat-type:4}; $6/8$ and
$12/8$ share \code{beat-type:8}).

\textbf{Bare accidental tokens.}
Accidentals appear as bare words (\code{flat}, \code{sharp},
\code{natural}) rather than prefixed (\code{acc:flat}).
This matches the OLIMPIC LMX
convention~\cite{mayerPracticalEndToEnd2024} and keeps the token
sequence compact.
Accidentals appear \textit{after} the duration in the token stream
because that mirrors the temporal order in which a reader processes a
note: pitch position, then duration, then any required display marks.

\textbf{Display-only accidentals.}
An accidental token (\code{flat}, \code{sharp}, \code{natural}) signals
that a flat, sharp, or natural sign should be \textit{drawn} on the
staff; it does not encode the pitch's chromatic alteration.
The chromatic pitch is already determined by the pitch letter, the
octave, and the active key signature.
This is the rule used by MusicXML, and it makes the model's prediction
task more local: deciding whether to display an accidental depends only
on the current measure and the key signature, not on a global pitch
inference.
A consequence is that the converter must compute---per
note---whether the pitch's alteration conflicts with the active key
signature or with a prior alteration in the same measure, and emit an
accidental token only when a sign should be drawn.

\subsection{Vocabulary}

The model treats LMX tokens as opaque class labels.
Three special tokens are reserved by convention: index 0 is the CTC
blank, index 1 is the padding token used during batch collation, and
index 2 is the out-of-vocabulary token.
The remaining indices, sorted alphabetically, encode the LMX tokens
that appear in the training corpus.

The vocabulary covers seven categories: pitch letters (\code{pitch:A}
through \code{pitch:G}); octave numbers (\code{octave:0} through
\code{octave:8}); the eight duration words; the dot token; the three
accidentals; the \code{measure} and \code{rest} structural tokens; the
tie tokens; the clef token \code{clef:G2}; key tokens
\code{key:fifths:-6} through \code{key:fifths:7}; and the time-signature
tokens.
The full pitch and octave ranges are always reserved even if absent in
the corpus, so that out-of-distribution pitch predictions decode
gracefully rather than collapsing to \code{<unk>}.

\subsection{Why LMX over alternatives}

Section~\ref{sec:soa-representations} surveyed the available symbolic
output formats; this thesis adopts LMX rather than MusicXML directly,
the PrIMuS native formats, or \code{**kern}.
LMX is round-trippable to MusicXML---the format downstream music
applications expect---without requiring the model to emit verbose XML
tags; it is the same representation used by the
OLIMPIC reference work~\cite{mayerPracticalEndToEnd2024}, keeping the
present results directly comparable; and its flat, whitespace-separated
form is directly compatible with CTC training, with no
linearisation logic required at the model boundary.

% ─────────────────────────────────────────────────────────────────────────────
\section{Model Design}
\label{sec:design-model}

The melody recognition stage is a Convolutional Recurrent Neural
Network trained with the Connectionist Temporal Classification loss
(CRNN--CTC).
This section justifies that architectural choice and its principal
hyperparameters in conceptual terms.
The implementation details---PyTorch module structure, exact stride
table, optimiser state---live in Chapter~\ref{ch:impl}.

\begin{figure}[ht]
  \centering
  \fbox{\parbox[c][6cm][c]{0.8\textwidth}{\centering \todo{FIGURE: CRNN--CTC block
  diagram. Input staff image $(B,1,128,W)$ $\to$ CNN backbone $\to$
  feature map $(B,C,1,W/4)$ $\to$ squeeze-height + transpose $\to$
  BiLSTM $\times$ 2 $\to$ linear projection to vocabulary $\to$
  log-softmax $\to$ CTC.}}}
  \caption{Block diagram of the CRNN--CTC melody recogniser.
  The CNN compresses height to one while preserving width up to a $4\times$
  factor; the resulting feature sequence is consumed left-to-right by
  the BiLSTM and decoded by CTC.}
  \label{fig:crnn-arch}
\end{figure}

\subsection{Why CRNN--CTC}

A single staff from a jazz lead sheet is read strictly left to right and
produces a variable-length sequence of tokens that is shorter than the
image is wide.
This matches the data assumptions of CTC\todo{ref? Graves et al.\ 2006
``Connectionist Temporal Classification'' --- not in references.bib} exactly:
the loss learns the alignment between an input feature sequence and a
shorter output token sequence without per-frame supervision, and it
tolerates the two-dimensional context (relative vertical position on
the staff) needed to disambiguate pitch from notehead shape.
End-to-end CRNN--CTC has been the dominant architecture for monophonic
OMR since \textcite{calvoZaragozaEndToEndNeural2018} and remains the
strongest baseline for the present
domain~\cite{calvoZaragozaCameraPrimus2018,mayerPracticalEndToEnd2024}.
A convolutional backbone provides translation-invariant feature
extraction along the staff width; a bidirectional recurrent layer
captures the local dependencies needed to interpret accidentals, dots,
and beam groupings; and CTC blank-collapsing decode produces a clean
symbolic output.
A transformer of comparable capacity would in principle perform at
least as well but at substantially higher data and compute costs that
are not justified for monophonic input
(Section~\ref{sec:soa-transformer}).

\subsection{Backbone choice}

The default backbone is a modified
ResNet18\todo{ref? He et al.\ 2016 ``Deep Residual Learning'' --- not in
references.bib} with asymmetric strides that compress the height
dimension aggressively while preserving width, leaving the horizontal
dimension as the CTC time axis.
A five-block VGG-style stack is retained as a legacy alternative for
backward compatibility with earlier checkpoints; it converges more
slowly and is not used for new experiments.
ResNet18 was preferred for three reasons: its residual connections give
the gradient a direct path through the network and accelerate
convergence; its pre-trained weights from ImageNet provide a reasonable
initialisation for the first convolutional layers, even though the
downstream task is monochrome rather than colour; and its parameter
count fits comfortably within the VRAM budget of a single consumer GPU
(see Chapter~\ref{ch:mgmt}).

\subsection{Temporal resolution and width compression}

The backbone compresses the width of the input image by a factor of
four, leaving the model with one feature vector per four input columns
of the original image.
This factor was chosen by a trade-off: a smaller compression ratio
preserves more time steps and so admits longer output sequences without
saturating CTC alignment, but increases compute cost; a larger ratio
risks producing more output tokens than time steps, which CTC's
\code{zero\_infinity} option silently masks but which corrupts gradient
estimates.
At a factor of four, all in-scope samples comfortably satisfy the CTC
length constraint, and the time axis remains rich enough to disambiguate
adjacent notes even on densely-engraved staves.

\subsection{Recurrent layer}

The convolutional features are fed to a stacked bidirectional
LSTM\todo{ref? Hochreiter \& Schmidhuber 1997 ``Long Short-Term Memory''
--- not in references.bib}.
Two layers are used, with 256 hidden units per direction; this gives
the model bidirectional context over the entire staff while keeping
parameter count modest.
A single-layer LSTM was found to under-fit on dense passages
(see ablations in Chapter~\ref{ch:results}); a three-layer LSTM gave no
measurable gain over two layers and slowed training.
Dropout between LSTM layers, applied during training only, regularises
the model against the inevitable overfitting that arises from training
on a synthetic corpus whose support is narrower than the inference-time
distribution.

% ─────────────────────────────────────────────────────────────────────────────
\section{Two-Stream Architecture for Lead Sheets}
\label{sec:design-twostream}

A jazz lead sheet contains two parallel visual streams: the melody on
the staff and the chord symbols typeset above it.
Their visual nature is fundamentally different.
The melody is a sequence of graphical symbols positioned relative to the
staff lines; correct interpretation depends on the clef and key
signature.
The chord symbols are short alphanumeric strings (\code{Cmaj7},
\code{F\#m7b5}, \code{Bb13}); their interpretation is a text-recognition
task and does not depend on staff geometry.

The architectural response is a \textit{two-stream} design: two
independent neural networks, each specialised to its own visual
distribution, with results merged in post-processing.
The melody network is the CRNN--CTC described in
Section~\ref{sec:design-model}; its vocabulary is the LMX token set.
The chord network is a separate CRNN--CTC of the same family but
character-level: its vocabulary is the 26-character alphabet sufficient
to spell every chord symbol in the Real Book.

\subsection{Why two networks rather than one}

A single network trained to emit both LMX tokens and chord characters
would have to learn two unrelated visual distributions and a vocabulary
that is the union of their token sets, with most tokens absent from any
given input region.
The two-stream design instead lets each network specialise: its training
data is purely in-domain, its vocabulary is minimal, and its training
schedule can be set independently.
Both surveys of the field~\cite{calvoZaragozaUnderstandingOMR2020,shatriOPTICALMUSICRECOGNITION}
explicitly recommend this decomposition for systems that handle
text-like adjuncts to staff notation.

\subsection{Independent training, joint inference}

The two networks are trained on disjoint datasets and produce
independent outputs.
At inference time, the staff detection stage emits both a staff image
and a chord-strip image for each detected system; the two are processed
in parallel and their outputs are joined into a single JSON segment
keyed by the system's bounding box on the page.
No cross-attention or shared representation links the two networks;
this is a deliberate simplicity choice.

\subsection{Chord vocabulary and notation conventions}

Chord symbols in the Real Book follow a stylised convention that
differs from many notation packages: minor chords are written with a
hyphen (\code{G-7} rather than \code{Gm7}), major-seventh chords with
\code{maj} (\code{Fmaj7}), half-diminished with the
\code{ø} glyph (\code{Bø}), and diminished with the spelled-out
\code{dim} suffix (\code{Cdim7}).
The character vocabulary covers the seven root letters, two
accidentals, four quality letters that spell \code{maj}, \code{dim},
and \code{sus}, the symbols \code{-}, \code{+}, \code{ø} for minor,
augmented, and half-diminished, the digits used in extensions
($1$, $3$, $6$, $7$, $9$), and the separator characters (space and
slash).
Out-of-scope chord features---\code{alt}, double accidentals,
parenthesised alterations---are deliberately absent from the
vocabulary and are mapped to nearest in-vocabulary spellings during
label canonicalisation.

% ─────────────────────────────────────────────────────────────────────────────
\section{Training Strategy}
\label{sec:design-training}

\subsection{Two-phase plan}

The training procedure has two distinct phases.
\textit{Phase 1} trains the model from scratch on the synthetic
LilyJAZZ corpus, with both clean and augmented versions of every
sample included in each epoch.
\textit{Phase 2} fine-tunes the resulting checkpoint on a small
manually-annotated set of Real Book pages, injected into the training
split only.
The validation and test splits remain drawn exclusively from the
synthetic corpus throughout, so that phase-2 fine-tuning does not
contaminate the held-out evaluation.

This is a textbook synthetic-pretrain → real-fine-tune transfer
learning plan; its effectiveness is the principal experimental
question of Chapter~\ref{ch:results}.

\subsection{Loss and supervision signal}

The supervision signal is the LMX token sequence associated with each
image.
CTC loss takes the log-probability tensor (one vector per output time
step), the target token sequence, and their respective lengths, and
computes the negative log-likelihood of the target marginalised over
all valid alignments to the input time axis.
The CTC blank class (index 0) absorbs time steps where no symbol is
emitted, which is the mechanism by which the model handles
variable-rate output without explicit alignment.

\subsection{Data splits}

The corpus is partitioned deterministically into train, validation, and
test splits with a fixed seed.
Validation and test fractions are each \SI{10}{\percent}; the remaining
\SI{80}{\percent} is used for training.
The same sample identifier always falls into the same split across
runs, so successive experiments are directly comparable.
When fine-tuning data is added, it is injected only into the train
split; no fine-tuning sample ever enters validation or test.

\subsection{Header-stripping augmentation}

Real Book pages often span multiple staves per system; every staff
after the first omits the clef, key, and time-signature header.
To prevent the model from over-relying on the header, a fraction of
training samples (\SI{40}{\percent}) have their header pixels cropped
out and the corresponding tokens removed from the label.
This forces the model to learn pitch and rhythm in the absence of an
explicit clef, which it can do because the staff geometry alone is
sufficient given the treble-clef-only contract.

\subsection{Rare-token oversampling}

Two LMX tokens (\code{tied:start} and \code{tied:stop}) are visually
subtle and statistically rare in the corpus, so the model
under-predicts them.
The training loader duplicates---in the training split only---every
sample whose label contains a rare token, with a configurable
multiplier.
This is a class-imbalance fix applied at the sample level rather than
in the loss function, which CTC does not easily admit.

% ─────────────────────────────────────────────────────────────────────────────
\section{Grammar Correction as a Design Layer}
\label{sec:design-grammar}

The CRNN model is trained to maximise the probability of the correct
token sequence; nothing in the training objective forces its output to
be a syntactically valid LMX sequence.
On most samples it is, but on dense or degraded inputs the raw
prediction can contain malformed fragments: a \code{pitch:C} not
followed by an \code{octave:} token, an accidental before the duration
of its note, a barline that arrives before the accumulated beats fill
the current measure.

A dedicated \term{grammar fixer} sits between the model and the
JSON output.
Its responsibilities are:

\begin{itemize}
    \item enforce the LMX grammar laid out in
          Section~\ref{sec:lmx-format} (correct token order within a
          note, correct duration after \code{rest}, correct nesting of
          tied groups);
    \item normalise predictions outside the lead-sheet contract: any
          clef other than \code{clef:G2} is rewritten to \code{clef:G2};
          time signatures outside the allowed eight-element set are
          rejected;
    \item clamp predicted octaves to the in-domain range $[3, 7]$,
          preventing rare misreads from producing unreachable pitches;
    \item count cumulative beats per measure against the active time
          signature and insert implicit barlines when the count
          overflows, producing well-formed measure boundaries even when
          the model misses a printed barline.
\end{itemize}

The grammar fixer is purely deterministic: it never invents new musical
content, only repairs structural defects in the model's output.
A learned alternative---training the model end-to-end against a
grammar-aware loss---is conceivable but adds complexity and offers
limited gain for a vocabulary as small and as well-structured as LMX.
Keeping grammar enforcement out of the loss also keeps the model's
output legible: every token in the raw prediction is interpretable in
the same way as the corresponding ground-truth token.

% ─────────────────────────────────────────────────────────────────────────────
\section{Evaluation Design}
\label{sec:design-eval}

\subsection{Primary metric: Symbol Error Rate}

The primary metric is the Symbol Error Rate~\cite{SymbolErrorRate}:
the Levenshtein edit distance between the predicted token sequence and
the ground truth, normalised by the length of the ground truth.
SER is reported on both the clean and the scanned versions of the test
split, so that the contribution of the augmentation pipeline is
visible.

\subsection{Category-level error breakdown}

SER aggregates errors of very different musical significance: a missed
barline counts identically to a wrong pitch.
To make this distribution legible, the evaluation reports a
category-level breakdown of edits attributable to each LMX token
family: \code{measure} tokens, tie tokens, pitch letters, octaves,
durations, accidentals, and key/time-signature tokens.
This is the same diagnostic information that OMR-NED provides as a
single number~\cite{martinezSevillaOMRNED2025}, expressed here as a
table rather than committed to a specific normalisation.

\subsection{Melodic SER}

A complementary \textit{melodic} SER is also reported, computed on the
same edit distance after stripping \code{measure} and tie tokens.
Structural tokens dominate the absolute edit count on the current
corpus but are largely musically irrelevant once a downstream renderer
re-bars the output; the melodic SER captures how often the model gets
the \textit{actual notes} right.

\subsection{Domain-gap evaluation}

The principal experimental question of this thesis is whether
synthetic pretraining plus fine-tuning closes the gap between training
on rendered images and inference on real scans.
Two ablations support this evaluation:
(i)~the model trained only on clean LilyJAZZ renders is evaluated on
real Real Book pages, exposing the raw synthetic-to-real gap;
(ii)~the same model after fine-tuning on a small annotated real set is
evaluated on the same pages, quantifying the reduction in the gap.
The fine-tuning protocol is the same as that described in
Section~\ref{sec:design-training}.
